\documentclass{article}

\title{Partikeln im Deutschen}
\date{}

\begin{document}
    
    \maketitle
    
    
    \section{Partikeln im Deutschen}
        
        \subsection{Vergleichspartikel}
            
            \textbf{Definition:} Vergleichspartikeln werden benutzt, um zwei Dinge miteinander zu vergleichen.
            
            \textbf{Beispiele:}
            
            \begin{itemize}
                
                \item \textbf{wie}
                
                Er ist so groß \textbf{wie} sein Bruder.\\
                He is as tall as his brother.
                
                \item \textbf{als}
                
                Sie ist älter \textbf{als} ich.\\
                She is older than I am.
            
            \end{itemize}
        
        \subsection{Gradpartikel}
            
            \textbf{Definition:} Gradpartikeln zeigen den Grad oder die Intensität einer Eigenschaft.
            
            \textbf{Beispiele:}
            
            \begin{itemize}
                
                \item \textbf{sehr}
                
                Das Buch ist \textbf{sehr} interessant.\\
                The book is very interesting.
                
                \item \textbf{ziemlich}
                
                Der Film war \textbf{ziemlich} gut.\\
                The movie was quite good.
                
                \item \textbf{so}
                
                Das ist \textbf{so} schön.\\
                That is so beautiful.
            
            \end{itemize}
        
        \subsection{Modalpartikel}
            
            \textbf{Definition:} Modalpartikeln drücken die Haltung oder Einstellung des Sprechers aus.
            
            \textbf{Beispiele:}
            
            \begin{itemize}
                
                \item \textbf{doch}
                
                Komm \textbf{doch} mit!\\
                Do come along!
                
                \item \textbf{ja}
                
                Du bist \textbf{ja} schon hier.\\
                You are already here, as I can see.
                
                \item \textbf{halt}
                
                Das ist \textbf{halt} so.\\
                That is just the way it is.
            
            \end{itemize}
        
        \subsection{Fokuspartikel}
            
            \textbf{Definition:} Fokuspartikeln betonen oder beschränken einen bestimmten Teil des Satzes.
            
            \textbf{Beispiele:}
            
            \begin{itemize}
                
                \item \textbf{nur}
                
                Ich habe \textbf{nur} Kaffee gekauft.\\
                I bought only coffee.
                
                \item \textbf{auch}
                
                Er kommt \textbf{auch} morgen.\\
                He is also coming tomorrow.
                
                \item \textbf{sogar}
                
                Sie hat \textbf{sogar} das ganze Buch gelesen.\\
                She even read the whole book.
            
            \end{itemize}
        
        \subsection{Negationspartikel}
            
            \textbf{Definition:} Negationspartikeln werden verwendet, um eine Aussage zu verneinen.
            
            \textbf{Beispiele:}
            
            \begin{itemize}
                
                \item \textbf{nicht}
                
                Ich habe das \textbf{nicht} gesehen.\\
                I did not see that.
            
            \end{itemize}
        
        \subsection{Antwortpartikel}
            
            \textbf{Definition:} Antwortpartikeln werden für kurze Antworten benutzt.
            
            \textbf{Beispiele:}
            
            \begin{itemize}
                
                \item \textbf{ja}
                
                Kommst du morgen? -- \textbf{Ja}.\\
                Are you coming tomorrow? -- Yes.
                
                \item \textbf{nein}
                
                Hast du Hunger? -- \textbf{Nein}.\\
                Are you hungry? -- No.
                
                \item \textbf{doch}
                
                Hast du kein Geld? -- \textbf{Doch}.\\
                Don't you have money? -- Yes, I do.
            
            \end{itemize}
        
        \subsection{Gliederungspartikel}
            
            \textbf{Definition:} Gliederungspartikeln strukturieren Texte oder Argumente.
            
            \textbf{Beispiele:}
            
            \begin{itemize}
                
                \item \textbf{erstens}
                
                \textbf{Erstens} müssen wir das Problem verstehen.\\
                First, we must understand the problem.
                
                \item \textbf{zweitens}
                
                \textbf{Zweitens} brauchen wir eine Lösung.\\
                Second, we need a solution.
                
                \item \textbf{übrigens}
                
                \textbf{Übrigens}, ich habe ihn gestern gesehen.\\
                By the way, I saw him yesterday.
            
            \end{itemize}

\end{document}